\section{模型假设与符号说明}
\subsection{模型假设}

结合赛题附件给出的参数口径与山区应急救援的实际情形，本文作如下假设：

\begin{enumerate}[label=\textbf{假设 \arabic*}]
	\item \textbf{航段净空假设}：无人机在航段上的巡航海拔取该航段所经过 DEM 像元的\textbf{最高地面高程 $+50$\,m}，即认为沿线地形与障碍物均不超过该高度，从而保证全程安全净空。
	\item \textbf{作业高度假设}：调度中心 O01 的作业高度取其地面海拔；各服务区的作业高度取地面海拔 $+30$\,m。连续访问多个服务区时，每次投送后均\textbf{从 30\,m 作业高度重新爬升}。
	\item \textbf{下降能耗假设}：附件给出的“下降能耗效率”为 0，故本文\textbf{不单独计算下降附加能耗}，下降段仅计入时间。
	\item \textbf{巡航能耗口径假设}：附件仅给出运输机的空载/满载标准航程与电池可用能量，未给巡航功率。本文按“飞满一个标准航程恰好耗尽一个可用能量”这一自洽口径\textbf{由航程反推巡航能耗}：$E^{hor}=(d/L_{g}(q))E_{g}^{use}$。
	\item \textbf{爬升效率假设}：附件列名为“爬升能耗效率”（取 0.72），本文按\textbf{效率}使用，即 $E^{up}=mgh^{+}/\eta_{up}$。
	\item \textbf{返航安全假设}：每个架次完成全部投送后必须返回 O01，返航为\textbf{空载}；架次总能耗须满足 $E_{p}^{T}\le(1-\rho_{g})E_{g}^{use}$。
	\item \textbf{资源独立性假设}：运输共享电池与中继能源组件为相互独立的资源，各自记录荷电状态，初始 SOC 均为 100\%；\textbf{不同机型的电池不可混用}；同一资源的任务占用与充电时段不得重叠，不同资源可并行充电。
	\item \textbf{通信三态假设}：链路状态分为直连、中继与中断三态，优先级为\textbf{直连 $>$ 中继 $>$ 中断}；中继需接入段与回传段\textbf{同时}可用；\textbf{不允许多跳转发}，每架运输机在任一时刻只由 G01 或一架中继保障。
	\item \textbf{地形遮挡假设}：DEM 为含植被与建筑的数字表面模型（DSM），本文按原样使用，遮挡判定结果偏保守；遮挡仅作为附加路径损耗（10\,dB）计入链路预算，而不直接判为中断。
	\item \textbf{分区刚性假设}：问题四中，问题三已确定的货箱组批、服务区访问顺序、任务安排与通信保障关系保持不变，且\textbf{同一运输架次涉及的服务区必须划入同一任务组}，任务执行期间各类资源不得跨组调配。
\end{enumerate}

\subsection{符号说明}

本文使用的主要符号及其含义如表~\ref{tab:symbols} 所示，未列出的符号在首次出现处说明。

\begin{table}[htbp]
	\caption{主要符号说明}
	\label{tab:symbols}
	\centering
	\zihao{-5}
	\setlength{\tabcolsep}{4pt}
	\begin{tabular}{llll}
		\toprule
		符号 & 含义 & 符号 & 含义 \\
		\midrule
		$O01$ & 临时调度中心（与网关 G01 同址） & $S_{i}$ & 第 $i$ 个服务区，$i=1,\dots,15$ \\
		$g$ & 运输机型，$g\in\{A,B,C\}$ & $R_{k}$ & 第 $k$ 架中继无人机，$k=1,2$ \\
		$Q_{g}$ & 机型 $g$ 的最大载货质量（kg） & $V_{g}$ & 机型 $g$ 的可用装载体积（$m^{3}$） \\
		$L_{g0},\,L_{gF}$ & 空载 / 满载标准航程（m） & $L_{g}(q)$ & 载荷 $q$ 时的等效航程（m） \\
		$E_{g}^{use}$ & 机型 $g$ 单组电池可用能量（kWh） & $\rho_{g}$ & 返航安全余量比例（附件取 0.20） \\
		$v_{g}^{\uparrow},\,v_{g}^{c},\,v_{g}^{\downarrow}$ & 爬升 / 巡航 / 下降速度（m/s） & $\eta_{up}$ & 爬升能耗效率（取 0.72） \\
		$H_{cruise}$ & 航段计划巡航海拔（m） & $H_{op}(n)$ & 节点 $n$ 的作业高度（m） \\
		$h^{+}_{ij},\,h^{-}_{ij}$ & 航段爬升 / 下降高度（m） & $d_{ij}$ & 节点 $i,j$ 间水平距离（m） \\
		$t_{gij}$ & 航段飞行时间（s） & $E_{gij}(q)$ & 航段能耗（kWh） \\
		$t_{chg}(s)$ & 由 SOC $=s$ 充满所需时间（s） & $T_{full}$ & 等效完全充电时间（s） \\
		$m_{b},\,v_{b}$ & 货箱 $b$ 的质量（kg）/ 体积（$m^{3}$） & $q_{max}^{safe}$ & 最大安全载荷（kg） \\
		$p$ & 一个运输架次 & $E_{p}^{T}$ & 架次总能耗（kWh） \\
		$P_{t,a},\,G_{a}$ & 发射端 $a$ 的发射功率（dBm）/ 天线增益（dBi） & $P_{sens,b}$ & 接收端 $b$ 的接收灵敏度（dBm） \\
		$M_{b}$ & 接收端 $b$ 的衰落裕量（dB） & $P_{th,b}$ & 接收端 $b$ 的有效接收门限（dBm） \\
		$L_{max,a\to b}$ & $a\to b$ 方向最大允许路径损耗（dB） & $L_{max,a\leftrightarrow b}$ & 双向链路门限（dB） \\
		$L_{FSPL}$ & 自由空间路径损耗（dB） & $L_{obs}$ & 地形遮挡附加损耗（取 10\,dB） \\
		$b_{ijt}$ & 时刻 $t$ 链路 $(i,j)$ 的遮挡指示变量 & $A_{ijt}$ & 时刻 $t$ 链路可用性指示变量 \\
		$D_{ijt}$ & 时刻 $t$ 链路两端三维距离（km） & $f$ & 载波频率（MHz） \\
		$P_{hover},\,P_{comms}$ & 中继悬停 / 通信附加功率（kW） & $K$ & 任务组数，$K\in\{1,2,3\}$ \\
		\bottomrule
	\end{tabular}
\end{table}
